% formula_supplement.tex — Color-coded formula decoder for
% "Audit the Scaffold, Not the Checkpoint" (companion to paper.pdf).
% Standalone: does NOT \input paper.tex or paper_appendix.tex, and cannot affect the paper build.
% Build:  pdflatex formula_supplement.tex   (no bibtex; no citations)
\documentclass[11pt]{article}

\usepackage[margin=1in]{geometry}
\usepackage{amsmath,amsfonts,amssymb,amsthm}
\usepackage{xcolor}
\usepackage{booktabs}
\usepackage{enumitem}
\usepackage{hyperref}
\hypersetup{colorlinks=true,linkcolor=black,urlcolor=black}

% Shared headline-statistics macros, the same single source of truth paper.tex
% uses. Every empirical value quoted below MUST come from a \stat... macro
% rather than a literal: this file previously hard-coded its own copies and
% drifted badly from the paper (Part B K, and all six Part C statistics --
% including a "directional" tier gap the paper reports as a null).
% stats_macros.tex — single source of truth for headline statistics quoted
% more than once across the NeurIPS workshop main body and appendix.
% \input by paper.tex, which also \inputs paper_appendix.tex, so the main body and
% the appendix cannot drift from each other (or from themselves) on a number.
%
% Naming: \statPB... = Part B (SWE-bench ablation), \statPC... = Part C
% (production sessions). Values are transcribed from the source analysis
% (scripts/part_b_anova_power.py, scripts/part_b_paper_numbers.py,
% scripts/part_b_edge_overlap.py, and the Part C statistical outputs) at full
% stated precision; every in-text mention of these quantities should use the
% macro, not a hardcoded literal.
%
% The blocks that scripts/generate_stats_macros.py patches automatically are
% Part C Finding 1, the Part B Model effect, the three Part B blocks below marked
% "measured weak-learning edge", "edge CIs", and "failMult overlap", the block
% marked "decomposing gamma's complement", the two \statPB*Eta macros, and the
% Part C version/release counts. Do not hand-edit those values -- re-run the
% generating script instead. Anything NOT in that set is hand-transcribed and
% therefore outside --check, which is the failure mode this boundary invites:
% the \statPB*Eta pair was hand-maintained until 2026-08-14 and \statPBSonnetEta
% had drifted from every arm's real value, and \statPBEdgeN plus the two
% \statPB*Edge series were likewise transcribed while being cited in the main
% text until the edge-decomposition work brought them under the generator.
% Prefer generating a macro over transcribing it.

% ---- Part B: design ----
\newcommand{\statPBK}{55}
\newcommand{\statPBSeeds}{5}
\newcommand{\statPBRounds}{4}
\newcommand{\statPBTrajectories}{1{,}650}

% ---- Part B: 2x3 repeated-measures ANOVA ----
\newcommand{\statPBModelF}{53.46}
\newcommand{\statPBModelDf}{(1, 54)}
\newcommand{\statPBModelP}{{<}0.0001}
\newcommand{\statPBModelEta}{0.497}
% Same effect with every row contributing its running best instead of its raw final
% round. The published DV gives a solved row its back-filled 1.0 but an unsolved row
% whatever its last attempt scored, and the stronger tier solves more rows, so the
% asymmetry runs with the tier effect; this shows the effect does not rest on it.
\newcommand{\statPBModelEtaRunningBest}{0.514}
\newcommand{\statPBModelCohenF}{0.99}
\newcommand{\statPBModelCohenFCI}{$[0.738, \, 1.319]$}
\newcommand{\statPBModeF}{0.50}
\newcommand{\statPBModeDf}{(2, 108)}
\newcommand{\statPBModeP}{0.61}
\newcommand{\statPBModeEta}{0.009}
\newcommand{\statPBModeCohenF}{0.10}
\newcommand{\statPBInterF}{0.81}
\newcommand{\statPBInterDf}{(2, 108)}
\newcommand{\statPBInterP}{0.45}
\newcommand{\statPBInterEta}{0.015}
\newcommand{\statPBInterCohenF}{0.12}

% ---- Part B: solve rates (mean-over-seeds final-round quality >= tau=0.6) ----
\newcommand{\statPBHaikuIndep}{24/55 (44\%) $[31,57]$}
\newcommand{\statPBHaikuBlind}{26/55 (47\%) $[35,60]$}
\newcommand{\statPBHaikuDiag}{27/55 (49\%) $[36,62]$}
\newcommand{\statPBSonnetIndep}{47/55 (85\%) $[74,92]$}
\newcommand{\statPBSonnetBlind}{47/55 (85\%) $[74,92]$}
\newcommand{\statPBSonnetDiag}{49/55 (89\%) $[78,95]$}
\newcommand{\statPBHaikuSolveRange}{44--49\%}
\newcommand{\statPBSonnetSolveRange}{85--89\%}

% ---- Part B: power analysis ----
% Bootstrap power at the K=55 we ran, from data/part_b_k55_anova_power.json
% ["full_pool"]["bootstrap_power"]["55"]: B=0.10657, AB=0.16577. Both are
% ROUNDED (not truncated) to match the convention used everywhere else in this
% file and in scripts/generate_stats_macros.py -- AB was previously 16\%,
% i.e. truncated, while B was rounded, so the pair disagreed on convention.
\newcommand{\statPBModePower}{11\%}
\newcommand{\statPBInterPower}{17\%}
% Required K for 80% power, interpolated on the logit scale between the
% bracketing points of that same grid (mode: K=430/500; interaction: 300/400).
\newcommand{\statPBModeReqK}{467}
\newcommand{\statPBInterReqK}{364}
% Hand-transcribed, so state the arithmetic: 364/55 = 6.6 (interaction) and
% 467/55 = 8.5 (mode). "7--8x" rounded the mode end DOWN, which understated how
% underpowered the null is -- the wrong direction for a disclosure.
\newcommand{\statPBPowerMultiple}{7--9$\times$}

% ---- Part B: noise floor ----
\newcommand{\statPBSeedSD}{0.096}
\newcommand{\statPBSignalOverNoise}{1.7$\times$}
\newcommand{\statPBBetweenModeSignal}{0.057}

% ---- Part B: measured weak-learning edge (diagnostic arm, N=275/tier) ----
% GENERATED as of the edge-decomposition work. These three were hand-transcribed
% while being cited in the main text, i.e. outside --check, which is exactly the
% drift-prone shape the header warns about; they are now derived from the same
% cache as the CIs and eta series below.
\newcommand{\statPBEdgeN}{275}
\newcommand{\statPBHaikuEdge}{(0, -0.39, -0.46, -0.45)}
\newcommand{\statPBSonnetEdge}{(0, -0.33, -0.41, -0.46)}
% Per-round improvement eta_bar_t on the running best (Thm. refinement), rounds
% 1--3; eta_0 is 0 by construction. GENERATED -- computed from the same pooled
% diagnostic-arm matrix as the gamma_t series above, so the two are guaranteed
% commensurable. Self-delimiting ($...$), so they are safe to cite inline.
\newcommand{\statPBHaikuEta}{$0.097 \to 0.036 \to 0.033$}
\newcommand{\statPBSonnetEta}{$0.156 \to 0.068 \to 0.029$}

% ---- Part B: edge CIs (task-cluster bootstrap over the 55 tasks) ----
% Rounds 1--3 only; round 0's gamma is fixed to 0 (no predecessor), so its
% interval is the degenerate [0,0] and is not quoted.
\newcommand{\statPBEdgeCIReps}{9{,}999}
\newcommand{\statPBHaikuEdgeCI}{$[{-}0.449, \, {-}0.322]$, $[{-}0.485, \, {-}0.431]$, $[{-}0.478, \, {-}0.416]$}
\newcommand{\statPBSonnetEdgeCI}{$[{-}0.373, \, {-}0.285]$, $[{-}0.449, \, {-}0.365]$, $[{-}0.482, \, {-}0.435]$}

% ---- Part B: decomposing gamma's complement (GENERATED) ----
% gamma uses a strict `>`, so "not improved" pools exact ties with genuine
% regressions. Splitting them is what stops gamma's sign from being read as
% AdaBoost's eps > 1/2: a converged, never-regressing refiner scores the floor
% -0.5 on ties alone. Ranges are over rounds 1--3 and both tiers.
\newcommand{\statPBTieRange}{80--95\%}
\newcommand{\statPBRegressMax}{2.9\%}
\newcommand{\statPBHaikuFlat}{80.7\%}
\newcommand{\statPBSonnetFlat}{69.1\%}
% The back-filling that makes the pooled gamma above uninterpretable. The
% experiment stops calling the model on the round a task first passes and writes a
% hardcoded 1.0 into every later round; gamma charges each such cell as a tie, so
% the pooled estimator is dragged down in proportion to the solve rate -- the
% stronger tier pads more and posts a *more* negative pooled edge.
\newcommand{\statPBHaikuPadShare}{27.6\%}
\newcommand{\statPBSonnetPadShare}{53.9\%}
\newcommand{\statPBHaikuPadTieRange}{33--45\%}
\newcommand{\statPBSonnetPadTieRange}{74--86\%}
% Tie share net of the padding -- the genuine no-op rate, which is what supports
% "refinement turns idempotent". Ranges over rounds 1--3.
\newcommand{\statPBHaikuGenuineTieRange}{52--59\%}
\newcommand{\statPBSonnetGenuineTieRange}{13--21\%}
% Schapire-Singer confidence-rated bound on measured cells: an unchanged round is
% an abstention, contributing Z_t = W_0 + 2 sqrt(W_+ W_-) < 1 rather than counting
% as an error. Product over rounds 1--3; the binary formulation's clipped-gamma
% curve pins at 1.0 instead, which is an artifact of that formulation.
\newcommand{\statPBHaikuAbstainZ}{0.81}
\newcommand{\statPBSonnetAbstainZ}{0.67}
% Edge conditional on the round having changed anything at all. Reported in
% app:edge-decomposition only: strongly positive, but the denominator is small and
% selected on outcome, so it cannot carry the primary claim.
\newcommand{\statPBHaikuMoversEdge}{{+}0.233 to {+}0.500}
\newcommand{\statPBSonnetMoversEdge}{{+}0.346 to {+}0.393}
% Edge re-measured on rows where a model call actually happened (q_{t-1} < 1.0) --
% the estimator the paper reports. Excluding the back-filled cells above is not
% conditioning on an outcome but declining to average over non-events. The
% denominator shrinks with t, so the intervals are wider and are quoted alongside
% the point estimates.
\newcommand{\statPBHaikuEdgeHeadroom}{({-}0.345, {-}0.435, {-}0.411)}
\newcommand{\statPBSonnetEdgeHeadroom}{({-}0.084, {-}0.143, {-}0.276)}
\newcommand{\statPBHaikuEdgeHeadroomCI}{$[{-}0.426, \, {-}0.247]$, $[{-}0.474, \, {-}0.383]$, $[{-}0.460, \, {-}0.348]$}
\newcommand{\statPBSonnetEdgeHeadroomCI}{$[{-}0.173, \, {+}0.021]$, $[{-}0.286, \, {+}0.028]$, $[{-}0.389, \, {-}0.125]$}
\newcommand{\statPBHaikuHeadroomN}{194$\to$158 of 275}
\newcommand{\statPBSonnetHeadroomN}{113$\to$49 of 275}
% The final-round eligible count. Only the Sonnet value is cited (its round-3
% denominator is the smallest of the six and the appendix flags it as the least
% reliable cell); the Haiku twin is generated anyway to keep the per-tier loop
% symmetric, and is deliberately unused -- check_macro_usage reports it, by design.
% Finding it a home is not the fix: 158 already reaches the reader inside
% \statPBHaikuHeadroomN ("194->158 of 275"), so a second mention would be redundant
% prose, and deleting the declaration would make patch_macros raise on the key the
% generator still emits.
\newcommand{\statPBHaikuHeadroomNLast}{158}
\newcommand{\statPBSonnetHeadroomNLast}{49}
% Per-round series behind the ranges above, quoted in app:edge-decomposition.
\newcommand{\statPBHaikuTieSeries}{(0.887, 0.945, 0.949)}
\newcommand{\statPBSonnetTieSeries}{(0.800, 0.898, 0.953)}
\newcommand{\statPBHaikuRegressSeries}{(0.004, 0.015, 0.000)}
\newcommand{\statPBSonnetRegressSeries}{(0.029, 0.011, 0.007)}
\newcommand{\statPBHaikuRegressRows}{5}
\newcommand{\statPBSonnetRegressRows}{12}

% ---- Part B: failMult overlap estimator (Prop. diversity's m* < k/2) ----
% 30 workers = 6 arms (2 tiers x 3 feedback modes) x 5 seeds, on the shared
% 55-task pool; a worker "fails" a task when its running-best quality stays
% below tau. Same task-cluster bootstrap as the edge CIs above.
\newcommand{\statPBOverlapK}{30}
\newcommand{\statPBOverlapHalfK}{15}
\newcommand{\statPBOverlapMStar}{30}
\newcommand{\statPBOverlapViolate}{23/55 (42\%)}
\newcommand{\statPBOverlapVoteErr}{0.418}
\newcommand{\statPBOverlapVoteErrCI}{$[0.291, \, 0.545]$}
\newcommand{\statPBOverlapEpsBar}{0.356}
\newcommand{\statPBOverlapEpsBarCI}{$[0.277, \, 0.436]$}
\newcommand{\statPBOverlapMarkov}{0.712}
\newcommand{\statPBOverlapGap}{{+}0.062}
\newcommand{\statPBOverlapGapCI}{$[{-}0.018, \, {+}0.142]$}
\newcommand{\statPBOverlapGapPosPct}{94\%}
% The paired gap's sign is round-dependent, so the interval above is a fact about the
% final round and not about every round. LastExclRound is the largest N for which every
% round 0..N excludes zero, which is what makes "through round N" true by construction
% rather than by assertion; the other two are that round's gap and interval.
\newcommand{\statPBOverlapGapLastExclRound}{2}
\newcommand{\statPBOverlapGapLastExcl}{{+}0.089}
\newcommand{\statPBOverlapGapLastExclCI}{$[{+}0.011, \, {+}0.167]$}
\newcommand{\statPBOverlapArms}{6}
\newcommand{\statPBOverlapArmMStar}{6}
\newcommand{\statPBOverlapArmViolate}{28/55}

% ---- Part B: the same overlap within one capability tier (GENERATED) ----
% The k=30 pool mixes tiers, so the two halves of the result above come apart here.
% m* is a max over tasks and cannot fall under restriction: a task defeating all 30
% workers defeats the 15 of either tier, so m*=k survives and the headline is not an
% artifact of pooling a strong tier with a weak one. The gap against eps_bar does NOT
% survive -- within a tier it is ~0 with an interval straddling zero -- so the pooled
% +0.062 is a tier-pooling effect and the paper claims nothing from its sign.
\newcommand{\statPBOverlapTierK}{15}
\newcommand{\statPBOverlapSonnetMStar}{15}
\newcommand{\statPBOverlapSonnetViolate}{8/55}
\newcommand{\statPBOverlapSonnetVoteErr}{0.145}
\newcommand{\statPBOverlapSonnetEpsBar}{0.147}
\newcommand{\statPBOverlapSonnetGap}{{-}0.001}
\newcommand{\statPBOverlapSonnetGapCI}{$[{-}0.039, \, {+}0.038]$}
\newcommand{\statPBOverlapHaikuMStar}{15}
\newcommand{\statPBOverlapHaikuViolate}{30/55}
\newcommand{\statPBOverlapHaikuVoteErr}{0.545}
\newcommand{\statPBOverlapHaikuEpsBar}{0.565}
\newcommand{\statPBOverlapHaikuGap}{{-}0.019}
\newcommand{\statPBOverlapHaikuGapCI}{$[{-}0.057, \, {+}0.017]$}

% ---- Part B: the SWE-bench resolved bar (GENERATED) ----
% TAU=0.6 thresholds *fractional* FAIL_TO_PASS/PASS_TO_PASS passing, so the solve
% rates above are NOT comparable to published SWE-bench resolve rates -- a task with
% 60% of its tests passing counts as solved above and as unresolved by SWE-bench's own
% criterion. These are the same aggregation (seed-mean of final-round quality,
% thresholded per task) re-run at 1.0. The edge gamma_t is tau-free, so nothing in the
% edge series or the ANOVA moves with this threshold; only solve rates and the overlap.
% Note statPBHaikuResolvedIndep > statPBHaikuResolvedDiag: at the resolved bar the weak
% tier's mode ordering *reverses* relative to tau=0.6, which is why the point-estimate
% crossover block below is quoted as a hypothesis rather than a result.
\newcommand{\statPBTauResolved}{1.0}
\newcommand{\statPBHaikuResolvedRange}{29--38\%}
\newcommand{\statPBSonnetResolvedRange}{71--75\%}
\newcommand{\statPBHaikuResolvedIndep}{21/55 (38\%)}
\newcommand{\statPBHaikuResolvedDiag}{16/55 (29\%)}
\newcommand{\statPBSonnetResolvedDiag}{41/55 (75\%)}
\newcommand{\statPBOverlapMStarResolved}{30}
\newcommand{\statPBOverlapViolateResolved}{23/55}

% ---- Part B: point-estimate crossover (not significant; quoted as hypothesis) ----
\newcommand{\statPBHaikuIndepQ}{0.432}
\newcommand{\statPBHaikuBlindQ}{0.449}
\newcommand{\statPBHaikuDiagQ}{0.466}
\newcommand{\statPBSonnetIndepQ}{0.875}
\newcommand{\statPBSonnetBlindQ}{0.856}
\newcommand{\statPBSonnetDiagQ}{0.871}
\newcommand{\statPBContrastFRange}{0.56--2.16}
\newcommand{\statPBContrastPRange}{0.15--0.46}
\newcommand{\statPBInteractionPRange}{0.13--0.33}

% ---- Part B: harness-verification audit ----
\newcommand{\statPBAuditFailN}{23/50 (46\%)}
\newcommand{\statPBAuditRootCauses}{5}
\newcommand{\statPBTaskPoolRepos}{10}
\newcommand{\statPBTaskPoolRetained}{40}
\newcommand{\statPBTaskPoolBackfilled}{15}

% ---- Part C: dataset ----
% Re-derived 2026-08-12 on an expanded live repo survey (was 378/1,172/21/28).
\newcommand{\statPCSessions}{401}
\newcommand{\statPCCommits}{1{,}211}
\newcommand{\statPCRepos}{23}
\newcommand{\statPCDevelopers}{20}

% ---- Part C: model-version heterogeneity within each tier label ----
% GENERATED. Derived from the commit `model` trailers in company_sessions.json by
% part_c_control_stats.py (["model_versions"]) and formatted by
% generate_stats_macros.py. company_simulations.model_cell does the parsing, so the
% "(1M context)" suffix collapses onto its base version -- the same model at a
% longer context window, not a separate capability level. Finding 2's tier contrast
% pools within each bin, so its within-bin capability spread is uncontrolled; these
% counts are what the paper cites as the size of that pooling.
\newcommand{\statPCOpusVersions}{five}
\newcommand{\statPCSonnetVersions}{three}
\newcommand{\statPCOpusTrailers}{eight}
\newcommand{\statPCSonnetTrailers}{four}
\newcommand{\statPCVersionSpan}{generations~4.5 through~5 in both families}

% ---- Part C: Finding 1 — churn convergence ----
\newcommand{\statPCDecayPct}{70.0\%}
\newcommand{\statPCMedianRatio}{0.49}
\newcommand{\statPCPermCI}{$[0.400, \, 0.793]$}
\newcommand{\statPCPermP}{{<}0.0001}
\newcommand{\statPCPermN}{70}

% ---- Part C: Finding 2 — model-capability gap (null) ----
\newcommand{\statPCTauOpus}{0.55}
\newcommand{\statPCTauSonnet}{0.55}
\newcommand{\statPCNOpus}{158}
\newcommand{\statPCNSonnet}{177}
\newcommand{\statPCGap}{{-}0.003}
\newcommand{\statPCGapCI}{$[{-}0.075, \, {+}0.069]$}
\newcommand{\statPCGapPermP}{0.93}
\newcommand{\statPCGapMWUP}{0.91}
\newcommand{\statPCEvaluePoint}{1.10}
\newcommand{\statPCEvalueCI}{1.00}
\newcommand{\statPCStratGap}{{+}0.002}
\newcommand{\statPCStratP}{0.96}
\newcommand{\statPCStrata}{4}
\newcommand{\statPCChurnRhoOpus}{0.80}
\newcommand{\statPCChurnRhoSonnet}{0.75}
\newcommand{\statPCChurnGap}{0.043}
\newcommand{\statPCChurnGapP}{0.54}
\newcommand{\statPCChurnGapN}{123}
\newcommand{\statPCMultilevelCoef}{{+}0.009}
\newcommand{\statPCMultilevelCI}{$[{-}0.070, \, {+}0.088]$}
\newcommand{\statPCMultilevelP}{0.82}
\newcommand{\statPCVarRepo}{0.132}
\newcommand{\statPCVarDev}{0.062}

% ---- Part C: Finding 3 — geometric churn-decay model ----
\newcommand{\statPCRho}{0.77}
\newcommand{\statPCRhoCI}{$[0.65, \, 0.92]$}
\newcommand{\statPCSigma}{2.39}
\newcommand{\statPCSigmaCI}{$[2.24, \, 2.54]$}
\newcommand{\statPCPiPure}{0.071}
\newcommand{\statPCTrainN}{137}
\newcommand{\statPCTestN}{60}
\newcommand{\statPCKSHeldout}{0.18}
\newcommand{\statPCKSHeldoutP}{0.045}
% Held-out KS is seed-sensitive with the expanded sample: across 8 train/test
% seeds p ranges [0.0003, 0.44], rejecting at alpha=0.05 in 6/8 splits — no
% longer the clean non-rejection reported on the smaller (378-session) sample.
\newcommand{\statPCKSHeldoutSeedRange}{$[0.0003, \, 0.44]$, rejecting in 6/8 seeds}
\newcommand{\statPCKSInsample}{0.21}
\newcommand{\statPCKSInsampleP}{{<}0.0001}

% ---- Part C: Finding 4 — session length ----
\newcommand{\statPCMedianLen}{1}
\newcommand{\statPCMeanLen}{3.0}
\newcommand{\statPCMaxLen}{45}
\newcommand{\statPCSingleCommitPct}{50.9\%}

% ---- Part C: robustness check — pre-AI-adoption human baseline ----
% Cutoff derived 2026-08-12 from the earliest AI-co-author-trailer commit
% found anywhere in the org's repos (2025-09-18), minus a 3-year safety
% margin. Window length (315 days) matches the AI-tier data's own
% observation span, ending at the cutoff. Repos are a disjoint, older cohort
% from the 23 AI-tier repos (excluded by construction), not a repo-matched
% comparison — see caveats in the surrounding text.
% Developer identity is heuristically deduplicated (git author-name/email
% aliasing is confirmed common in this org's history: GitHub-noreply
% handles, contractor accounts, work/personal email pairs — see
% scripts/resolve_author_aliases.py) rather than a validated headcount; the
% prose deliberately does not assert the raw number as exact.
\newcommand{\statPCCtrlSince}{2021-11-08}
\newcommand{\statPCCtrlUntil}{2022-09-19}
\newcommand{\statPCCtrlRepos}{122}
\newcommand{\statPCCtrlDevelopers}{62}
\newcommand{\statPCCtrlSessions}{9{,}395}
\newcommand{\statPCCtrlCommits}{17{,}694}
\newcommand{\statPCCtrlTau}{0.39}
\newcommand{\statPCCtrlRho}{0.86}
\newcommand{\statPCCtrlRhoCI}{$[0.82, \, 0.91]$}
% Session-length composition. The control cohort is far more singleton-heavy than
% the AI tier, so the survivor-ratio gap is also reported restricted to
% multi-commit sessions -- see the robustness-check paragraph.
\newcommand{\statPCCtrlSingleCommitPct}{63.9\%}
\newcommand{\statPCCtrlTauMulti}{0.435}
\newcommand{\statPCTauMulti}{0.593}
\newcommand{\statPCCtrlGapOpus}{{-}0.160}
\newcommand{\statPCCtrlGapOpusP}{{<}0.0002}
\newcommand{\statPCCtrlGapSonnet}{{-}0.163}
\newcommand{\statPCCtrlGapSonnetP}{{<}0.0002}
\newcommand{\statPCCtrlMLOpusCoef}{{+}0.145}
\newcommand{\statPCCtrlMLOpusP}{0.00082}
\newcommand{\statPCCtrlMLSonnetCoef}{{+}0.154}
\newcommand{\statPCCtrlMLSonnetP}{0.00026}

% ---- Part C: release footprint (Responsible Use Statement) ----
% Distinct people appearing anywhere in the four published datasets. Strictly
% more than \statPCCtrlDevelopers, which counts only the pre-AI control roster:
% the AI-era cohorts contain people the control window never saw. The privacy
% claim is about everyone in the release, so it must not reuse the control count.
\newcommand{\statPCReleasedDevelopers}{71}
\newcommand{\statPCReleasedDatasets}{four}

% ---- Part C: tier composition of the AI-tier cohort ----
% Opus and Sonnet do not exhaust the cohort. A small Haiku residue and a few
% sessions whose trailer does not resolve sit outside Finding 2's contrast, and
% the appendix's capability scale ranks the Haiku cell, so the main text names
% them rather than implying only two tiers are present.
\newcommand{\statPCHaikuSessions}{9}
\newcommand{\statPCUnresolvedSessions}{5}
\newcommand{\statPCTierExcluded}{14}
% Resamples behind every Part C permutation test. Quoted wherever a p-value sits at
% the two-sided resolution floor 2/(reps+1), so "p < 0.0002" reads as "no resample
% reached the observation" rather than as an estimate that happens to be small.
\newcommand{\statPCPermReps}{9{,}999}

% ---- Part C: within-era AI-vs-non-AI contrast ----
% The one design free of BOTH the person confound (same developers on both sides)
% and the era confound (same window, codebase, tooling). Split only on whether a
% commit carries a Claude trailer, so what remains uncontrolled is task selection.
% All values from data/part_c_control_stats.json ["within_era"].
\newcommand{\statPCWithinSince}{2025-09-18}
\newcommand{\statPCWithinUntil}{2026-12-31}
\newcommand{\statPCWithinAISessions}{487}
\newcommand{\statPCWithinAICommits}{1{,}493}
\newcommand{\statPCWithinNonAISessions}{1{,}874}
\newcommand{\statPCWithinNonAICommits}{4{,}914}
\newcommand{\statPCWithinDevelopers}{22}
\newcommand{\statPCWithinTauAI}{0.540}
\newcommand{\statPCWithinTauNonAI}{0.419}
% Paired within developer, at four commit floors. Eight tests in total (two
% statistics x four floors); the floor labels are the min_commits_per_side the
% pairing was run at. Renaming a floor breaks the macro lookup loudly rather than
% silently retargeting a row, since patch_macros raises on a missing macro.
\newcommand{\statPCWithinFloorB}{15}
\newcommand{\statPCWithinFloorC}{30}
\newcommand{\statPCWithinFloorD}{50}
\newcommand{\statPCWithinFZeroPairs}{21}
\newcommand{\statPCWithinFZeroHigher}{14}
\newcommand{\statPCWithinFZeroSignP}{0.19}
\newcommand{\statPCWithinFZeroWilcoxP}{0.022}
\newcommand{\statPCWithinFZeroMedian}{{+}0.127}
\newcommand{\statPCWithinFBPairs}{14}
\newcommand{\statPCWithinFBHigher}{10}
\newcommand{\statPCWithinFBSignP}{0.18}
\newcommand{\statPCWithinFBWilcoxP}{0.011}
\newcommand{\statPCWithinFBMedian}{{+}0.094}
\newcommand{\statPCWithinFCPairs}{9}
\newcommand{\statPCWithinFCHigher}{6}
\newcommand{\statPCWithinFCSignP}{0.51}
\newcommand{\statPCWithinFCWilcoxP}{0.098}
\newcommand{\statPCWithinFCMedian}{{+}0.061}
\newcommand{\statPCWithinFDPairs}{7}
\newcommand{\statPCWithinFDHigher}{5}
\newcommand{\statPCWithinFDSignP}{0.45}
\newcommand{\statPCWithinFDWilcoxP}{0.11}
\newcommand{\statPCWithinFDMedian}{{+}0.127}
% The signed spread at the >=15 floor. This is what Wilcoxon is reacting to and the
% sign test is not: quoting the significant p-value without it would hide which
% assumption carries the result.
\newcommand{\statPCWithinFBDeltaMin}{{-}0.078}
\newcommand{\statPCWithinFBDeltaMax}{{+}0.332}
% Best-measured pair (most commits on BOTH sides), so "the direction is not an
% artifact of thin data" is checkable rather than asserted.
\newcommand{\statPCWithinTopVolNonAI}{1{,}340}
\newcommand{\statPCWithinTopVolAI}{444}
\newcommand{\statPCWithinTopVolDelta}{{+}0.127}

% ---- Part C: pre-AI within-developer pairing ----
% Cuts AGAINST the between-group human-vs-AI gap and is reported that way. From
% ["paired_within_developer"]; floor 0 and the >=50 floor only.
\newcommand{\statPCPairedPairs}{12}
\newcommand{\statPCPairedHigher}{9}
\newcommand{\statPCPairedSignP}{0.15}
\newcommand{\statPCPairedWilcoxP}{0.034}
\newcommand{\statPCPairedMedian}{{+}0.117}
\newcommand{\statPCPairedFloorD}{50}
\newcommand{\statPCPairedFDPairs}{4}
\newcommand{\statPCPairedFDHigher}{2}
\newcommand{\statPCPairedFDSignP}{1}
\newcommand{\statPCPairedFDMedian}{{+}0.017}
% Best-measured pair, and the worst inclusion asymmetry. The control keeps all of a
% person's commits while the AI tier keeps only their Claude-co-authored ones, so a
% large ratio means low adoption, not collapsed output.
\newcommand{\statPCPairedTopVolCtrl}{322}
\newcommand{\statPCPairedTopVolAI}{257}
\newcommand{\statPCPairedTopVolDelta}{{+}0.320}
\newcommand{\statPCPairedMaxRatio}{33}

% ---- Part C: capability-ordered re-test of Finding 2 ----
% Finding 2's tier contrast pools model versions, which would attenuate a real gap.
% This re-tests it as a monotone trend over an ordered capability scale, under two
% orderings (each the other's sensitivity analysis). From ["capability"].
\newcommand{\statPCCapSessions}{343}
\newcommand{\statPCCapCells}{8}
\newcommand{\statPCCapSteps}{7}
% Cell composition. The cell total coincides with the pooled-version total quoted
% beside it (five Opus + three Sonnet versions, eight cells), so a reader reconstructs
% "the cells are the Opus and Sonnet versions" and is wrong twice: one Opus version is
% the plurality cell of no session, and Haiku contributes a cell the tier contrast
% never sees. Stated explicitly for that reason.
\newcommand{\statPCCapCellsOpus}{four}
\newcommand{\statPCCapCellsSonnet}{three}
\newcommand{\statPCCapCellsHaiku}{one}
\newcommand{\statPCCapAbsentOpusVersion}{5}
\newcommand{\statPCCapFamilyRho}{{+}0.012}
\newcommand{\statPCCapFamilyP}{0.82}
\newcommand{\statPCCapFamilyCoef}{{+}0.003}
\newcommand{\statPCCapFamilyCoefP}{0.75}
\newcommand{\statPCCapFamilyCI}{$[{-}0.016, \, {+}0.022]$}
\newcommand{\statPCCapVersionRho}{{-}0.008}
\newcommand{\statPCCapVersionP}{0.88}
\newcommand{\statPCCapVersionCoef}{{-}0.011}
\newcommand{\statPCCapVersionCoefP}{0.38}
\newcommand{\statPCCapVersionCI}{$[{-}0.036, \, {+}0.014]$}
% The widest total swing the adjusted CI still permits across all rank steps. The
% null is BOUNDED, not tight: compare this against the human-vs-AI gap.
\newcommand{\statPCCapSwing}{0.155}
% Per-cell taus are not monotone in either ordering, which is why rho sits near
% zero. The Haiku cell rests on very few sessions.
\newcommand{\statPCCapTauSonnetMid}{0.569}
\newcommand{\statPCCapTauSonnetNew}{0.446}
\newcommand{\statPCCapTauHaiku}{0.352}
\newcommand{\statPCCapNHaiku}{8}


% --- Paper macros reused inside formulas (from paper.tex:38-44) ---
\newcommand{\R}{\mathbb{R}}
\newcommand{\X}{\mathcal{X}}
\newcommand{\Y}{\mathcal{Y}}
\newcommand{\D}{\mathcal{D}}
\newcommand{\Hyp}{\mathcal{H}}
\newcommand{\sign}{\mathrm{sign}}
\newcommand{\err}{\mathrm{err}}

% --- Semantic-role color palette (Paul Tol "muted", colorblind-safe) ---
\definecolor{specclr}{HTML}{332288}  % indigo  — specifications / inputs / distribution
\definecolor{tgtclr}{HTML}{117733}   % green   — targets / optimal output
\definecolor{agentclr}{HTML}{882255} % wine    — weak agent / hypothesis / operators
\definecolor{wtclr}{HTML}{AA4499}    % purple  — confidence weight / margin / normalizer
\definecolor{erfclr}{HTML}{CC6677}   % rose    — error / edge / loss / thresholds
\definecolor{valclr}{HTML}{44AA99}   % teal    — quality / value / improvement / ceiling
\definecolor{empclr}{HTML}{999933}   % olive   — empirical fit parameters

% One wrapper per role; works in math mode.
\newcommand{\spec}[1]{\textcolor{specclr}{#1}}
\newcommand{\tgt}[1]{\textcolor{tgtclr}{#1}}
\newcommand{\agent}[1]{\textcolor{agentclr}{#1}}
\newcommand{\wt}[1]{\textcolor{wtclr}{#1}}
\newcommand{\erf}[1]{\textcolor{erfclr}{#1}}
\newcommand{\val}[1]{\textcolor{valclr}{#1}}
\newcommand{\empc}[1]{\textcolor{empclr}{#1}}

% --- Legend helper: one row = colored term + plain-English explanation ---
% \leg{$colored term$}{explanation text}
\newcommand{\leg}[2]{\item[{\normalfont #1}] #2}
\newenvironment{legend}%
  {\par\smallskip\begin{description}[leftmargin=2.6em,style=nextline,font=\normalfont,itemsep=1pt,topsep=2pt]}%
  {\end{description}\par\smallskip}

\newcommand{\xref}[1]{\par\smallskip{\footnotesize\textit{Main text:} #1}}

\title{\textbf{Annotated Formula Supplement}\\[2pt]
{\large A Color-Coded Decoder for\\ \emph{Audit the Scaffold, Not the Checkpoint}}}
\author{}
\date{}

\begin{document}
\maketitle
\vspace{-3em}

\section*{How to read this supplement}
This companion document reproduces every defining and result formula from the main paper
and colors each term by its \textbf{semantic role}. A color means the same kind of thing
\emph{everywhere}: the confidence weight $\wt{\alpha_t}$ is purple in the protocol, in the
training-error proof, and in the algorithm alike. Structural operators
($\textstyle\sum,\ \prod,\ \ln,\ \exp$) and plain indices ($t,T,N,i,K$) are left black to
reduce clutter. Under each formula a short legend names its colored terms in plain English
and points back to where the formula appears in the main text.

\begin{center}
\renewcommand{\arraystretch}{1.35}
\begin{tabular}{@{}c l l@{}}
\toprule
\textbf{Color} & \textbf{Role} & \textbf{Typical symbols} \\
\midrule
$\spec{\blacksquare}$ & Specifications / inputs / spec-distribution & $\spec{\X},\ \spec{x_i},\ \spec{\D},\ \spec{D_t},\ \spec{w_i}$ \\
$\tgt{\blacksquare}$  & Artifacts: drafts and optimal targets & $\tgt{\Y},\ \tgt{y_i},\ \tgt{y},\ \tgt{z_t}$ \\
$\agent{\blacksquare}$ & Weak agent / hypothesis / operators & $\agent{h_t},\ \agent{H},\ \agent{\Hyp},\ \agent{R_t},\ \agent{F_T}$ \\
$\wt{\blacksquare}$   & Confidence weight / margin / normalizer & $\wt{\alpha_t},\ \wt{f(x)},\ \wt{Z_t},\ \wt{J}$ \\
$\erf{\blacksquare}$  & Error / edge / loss / thresholds & $\erf{\epsilon_t},\ \erf{\gamma_t},\ \erf{L},\ \erf{\theta},\ \erf{\delta}$ \\
$\val{\blacksquare}$  & Quality / value / improvement / ceiling & $\val{V},\ \val{V^*(W)},\ \val{\eta_t},\ \val{q_{i,t}},\ \val{B}$ \\
$\empc{\blacksquare}$ & Empirical fit parameters (Part~C) & $\empc{\rho},\ \empc{\sigma},\ \empc{\psi},\ \empc{\pi_{\mathrm{pure}}}$ \\
\bottomrule
\end{tabular}
\end{center}

\begin{quote}\small
\textbf{Two glyph warnings.}
\emph{(1)}~Two quantities that a single $\tau$ would conflate are kept distinct:
$\erf{\tau}=0.6$ is Part~B's \emph{pass threshold} (an error/edge quantity, rose), while
Part~C's \emph{mean survivor-ratio statistics} are $\empc{\psi_{\mathrm{Opus}}},
\empc{\psi_{\mathrm{Sonnet}}}$ (empirical, olive). Part~C's $\empc{\rho}$ is likewise the
churn-decay factor only; its rank correlations are written $r_s$.
\emph{(2)}~The success indicator flips: the theory writes $\erf{L}(\tgt{y},\agent{h}(\spec{x}))>0$
for \emph{failure}, while the algorithm's weight update writes $\erf{L}(\cdots)=0$ for
\emph{success} inside the exponent. Both encode the same AdaBoost update, read oppositely.
\end{quote}

\newpage
\section{Problem Setup}

\[
\erf{L}(\tgt{y}, \hat{\tgt{y}}), \qquad \agent{f} : \spec{\X} \to \tgt{\Y}, \qquad
\D \ \text{over}\ \spec{\X}\times\tgt{\Y}.
\]
\begin{legend}
\leg{$\erf{L}(\tgt{y},\hat{\tgt{y}})$}{structured loss between a candidate artifact $\hat{\tgt{y}}$ and the optimal target $\tgt{y}$ (e.g.\ token or AST edit distance).}
\leg{$\spec{\X}$}{space of specifications (prompts / issue tickets).}
\leg{$\tgt{\Y}$}{space of optimal code artifacts (the targets).}
\leg{$\spec{\D}$}{joint distribution over specification–solution pairs.}
\leg{$\agent{f}$}{the structured mapping specifications $\to$ code that the system learns.}
\end{legend}
\xref{Problem Setup, \S\,Theoretical Analysis.}

\section{Weak Coding Agent and Edge}

\[
\err_{\spec{\D}}(\agent{h_t}) = \Pr_{(\spec{x},\tgt{y})\sim\spec{\D}}\!\big[\erf{L}(\tgt{y},\agent{h_t}(\spec{x}))>0\big] \ \le\ \tfrac{1}{2}-\erf{\gamma_t}.
\]
\begin{legend}
\leg{$\agent{h_t}$}{the weak coding agent at round $t$: proposes a patch to improve the current draft.}
\leg{$\err_{\spec{\D}}(\agent{h_t})$}{its probability of failing (nonzero loss) on a spec drawn from $\spec{\D}$.}
\leg{$\erf{L}(\cdots)>0$}{the boolean failure criterion (loss $>0$ = fail; $=0$ = success).}
\leg{$\erf{\gamma_t}$}{the \emph{edge} — how much better than chance ($\tfrac12$) the agent is. In code space even a small positive edge presupposes real capability.}
\end{legend}
\xref{Definition~1 (Weak Coding Agent); edge revisited in Remark~14.}

\section{Strong Coding System (Two Levels)}

\[
\agent{F_T}(\spec{x}) = \big(\agent{R_T}\circ \agent{R_{T-1}}\circ\cdots\circ \agent{R_1}\big)\!\big(\agent{F_0}(\spec{x})\big)
\qquad\text{(artifact level)}
\]
\[
\wt{f}(\spec{x}) = \sum_{t=1}^{T} \wt{\alpha_t}\, \agent{h_t}(\spec{x}), \qquad \wt{\alpha_t}\in\R_{\ge 0}
\qquad\text{(score level: additive margin)}
\]
\begin{legend}
\leg{$\agent{F_T}$}{the artifact after $T$ rounds — a \emph{sequential composition} of patches (order-dependent, non-linear).}
\leg{$\agent{R_t}$}{refinement operator applying patch $\agent{h_t}$ to the current draft.}
\leg{$\agent{F_0}$}{the initial draft before any refinement.}
\leg{$\wt{f}(\spec{x})$}{the additive confidence \emph{margin} — a weighted vote over specifications; this is what inherits the AdaBoost machinery.}
\leg{$\wt{\alpha_t}$}{confidence weight assigned to agent $t$.}
\end{legend}
\xref{Definition~8 (Strong Coding System). Artifact vs.\ margin distinction (\S\,Framework).}

\section{Agentic Boosting Protocol}

\[
\erf{\epsilon_t} = \Pr_{(\spec{x},\tgt{y})\sim \spec{D_t}}\!\big[\erf{L}(\tgt{y},\agent{h_t}(\spec{x}))>0\big],
\qquad
\wt{\alpha_t} = \tfrac{1}{2}\ln\frac{1-\erf{\epsilon_t}}{\erf{\epsilon_t}},
\]
\[
\spec{D_{t+1}}(i) = \frac{\spec{D_t}(i)\,\exp\!\big(-\wt{\alpha_t}\,\tgt{y_i}\,\agent{h_t}(\spec{x_i})\big)}{\wt{Z_t}}.
\]
\begin{legend}
\leg{$\erf{\epsilon_t}$}{weighted error of agent $t$ under the current spec-distribution $\spec{D_t}$.}
\leg{$\wt{\alpha_t}$}{confidence weight — larger when the error is smaller (the AdaBoost rule).}
\leg{$\spec{D_t}(i)$}{weight on spec $i$ at round $t$; hard specs get up-weighted.}
\leg{$\wt{Z_t}$}{normalizer keeping $\spec{D_{t+1}}$ a probability distribution.}
\end{legend}
\xref{Definition~9 (Agentic Boosting Protocol).}

\section{Training-Error Bound}

\[
\frac{1}{N}\sum_{i=1}^{N}\mathbf{1}\!\big[\agent{H}(\spec{x_i})\neq \tgt{y_i}\big] \ \le\ \prod_{t=1}^{T}\wt{Z_t},
\qquad \wt{Z_t}=2\sqrt{\erf{\epsilon_t}\,(1-\erf{\epsilon_t})}.
\]
\begin{legend}
\leg{$\agent{H}$}{the strong system's decision (sign of the margin $\wt{f}$).}
\leg{$\mathbf{1}[\agent{H}\neq \tgt{y_i}]$}{training mistake indicator for spec $i$.}
\leg{$\prod_t \wt{Z_t}$}{product of per-round normalizers — the bound on training error.}
\leg{$\wt{Z_t}$}{shrinks below $1$ whenever the agent has positive edge, so the product decays.}
\end{legend}
\xref{Theorem~10 (Training Error of Agentic Boosting).}

\section{Exponential Decay of Training Error}

\[
\err_{\mathrm{train}}(\agent{H}) \ \le\ \exp\!\big(-2\,\erf{\gamma}^2\,T\big).
\]
\begin{legend}
\leg{$\err_{\mathrm{train}}(\agent{H})$}{fraction of training specs the strong system gets wrong.}
\leg{$\erf{\gamma}$}{uniform edge lower bound ($\erf{\gamma_t}\ge\erf{\gamma}>0$ for all $t$).}
\leg{$T$}{number of rounds; error falls exponentially in $T$ once the edge holds.}
\end{legend}
\xref{Corollary~11 (Exponential Decay). Applies only once agents cross the capability threshold (Remark~14).}

\section{Generalization Bound}

\[
\Pr_{(\spec{x},\tgt{y})\sim\spec{\D}}\!\big[\agent{H}(\spec{x})\neq \tgt{y}\big]
\ \le\ \Pr_{\mathrm{train}}\!\big[\tgt{y}\,\wt{f}(\spec{x})\le \erf{\theta}\big]
\ +\ O\!\left(\sqrt{\frac{\erf{d}\ln(N/\erf{d})\ln(1/\erf{\theta})+\ln(1/\erf{\delta})}{N}}\right).
\]
\begin{legend}
\leg{$\Pr_{\spec{\D}}[\cdots]$}{true (test) error over the data distribution.}
\leg{$\tgt{y}\,\wt{f}(\spec{x})\le\erf{\theta}$}{training margin-violation term; shrinks as boosting grows margins.}
\leg{$\erf{\theta}$}{margin threshold ($>0$).}
\leg{$\erf{d}$}{VC-dimension of the weak-agent class $\agent{\Hyp}$.}
\leg{$\erf{\delta}$}{failure probability (bound holds with probability $1-\erf{\delta}$).}
\leg{$N$}{number of i.i.d.\ training samples. The bound is \emph{independent of $T$}.}
\end{legend}
\xref{Theorem~12 (Generalization Error). Margin bound of Schapire et al.\ (1998).}

\section{Equivalence to Forward Stagewise Additive Modeling}

\[
\wt{J}(\wt{\alpha},\agent{h}) = \sum_{i=1}^{N}\exp\!\Big(-\tgt{y_i}\big(\wt{f_{t-1}}(\spec{x_i})+\wt{\alpha}\,\agent{h}(\spec{x_i})\big)\Big),
\qquad \erf{L}(\tgt{y},\wt{f}(\spec{x})) = e^{-\tgt{y}\,\wt{f}(\spec{x})}.
\]
\begin{legend}
\leg{$\wt{J}(\wt{\alpha},\agent{h})$}{exponential-loss objective minimized greedily each round.}
\leg{$\wt{f_{t-1}}$}{margin accumulated through round $t-1$.}
\leg{$\wt{\alpha}\,\agent{h}$}{the new term added this round (weight $\times$ agent).}
\leg{$\erf{L}=e^{-\tgt{y}\wt{f}}$}{the exponential loss whose greedy coordinate descent \emph{is} the protocol.}
\end{legend}
\xref{Theorem~13 (Boosting as Greedy Coordinate Descent).}

\section{The Refinement Game}

\[
\mathbb{E}_{\tgt{y}\sim\D(\cdot\mid \spec{x})}\big[\val{V}(\agent{R_t}(\tgt{z},a_t^*))\big] \ \ge\ \val{V}(\tgt{z}) + \val{\eta_t}
\qquad\text{(improvement property)}
\]
\[
\mathbb{E}\big[\val{V}(\tgt{z_T})\big] \ \ge\ \val{V}(\tgt{z_0}) + \sum_{t=1}^{T}\val{\eta_t},
\qquad
\sum_{t=1}^{\infty}\val{\eta_t} \ \le\ 1 - \val{V}(\tgt{z_0})
\quad\text{(saturation)}.
\]
\begin{legend}
\leg{$\val{V}$}{quality function $\tgt{\Y}\to[0,1]$ — how good the current artifact is. Note the domain: quality is scored on \emph{artifacts}, not on the specification space $\spec{\X}$.}
\leg{$\agent{R_t}(\tgt{z},a_t^*)$}{best refinement of draft $\tgt{z}$ under action $a_t^*$ (patch / refactor / test).}
\leg{$\val{\eta_t}$}{guaranteed minimum quality gain at round $t$ (analog of the edge $\erf{\gamma_t}$).}
\leg{$\tgt{z_0},\tgt{z_T}$}{initial and round-$T$ drafts, both artifacts in $\tgt{\Y}$.}
\leg{$\sum\val{\eta_t}\le 1-\val{V}(\tgt{z_0})$}{total improvement is capped: gains must vanish, $\val{\eta_t}\to 0$.}
\end{legend}
\xref{Definition~1 (Weak Coding Agent), which introduces $\agent{R_t}$, and Theorem~2 (Monotone Improvement under Refinement).}

\section{High-Probability Refinement}

\[
\Pr\!\left[\val{V}(\tgt{z_T}) \ge \val{V}(\tgt{z_0}) + \sum_{t=1}^{T}\val{\eta_t}\right] \ \ge\ 1 - \sum_{t=1}^{T}\erf{p_t}.
\]
\begin{legend}
\leg{$\erf{p_t}$}{per-round probability that the agent \emph{fails} to deliver the gain $\val{\eta_t}$ (a quality regression).}
\leg{$1-\sum\erf{p_t}$}{union-bound guarantee that all rounds improve, relaxing the strict expectation.}
\end{legend}
\xref{Corollary~15 (High-Probability Refinement). Verified against Mathlib's probability library.}

\section{Stationarity Dichotomy}

\[
\text{(i) fixed ceiling } \val{V}(\tgt{z})\le \val{B}:\quad
\sum_{t=1}^{\infty}\val{\eta_t} \le \val{B} - \mathbb{E}[\val{V}(\tgt{z_0})]
\ \Rightarrow\ \val{\eta_t}\to 0 \ \text{(saturation)}.
\]
\[
\text{(ii) sustained edge } \val{\eta_t}\ge \val{c}>0:\quad
\forall \val{B}\in\R,\ \exists\,T\ \text{s.t.}\ \mathbb{E}[\val{V}(\tgt{z_T})] > \val{B}
\ \text{(divergence)}.
\]
\begin{legend}
\leg{$\val{B}$}{a fixed upper bound (ceiling) on quality.}
\leg{$\val{c}$}{a uniform positive per-round gain the expanding class must sustain to diverge.}
\leg{(i)}{\emph{stationary} regime — fixed reachable class $\Rightarrow$ diminishing returns.}
\leg{(ii)}{\emph{non-stationary} regime — only a perpetually expanding class escapes the plateau.}
\end{legend}
\xref{Proposition~3 (Stationarity Dichotomy). Formally verified in Lean (\texttt{Proposition8.lean}).}

\section{The Frozen-Weight Ceiling}

\[
\text{(M1)}\ \ \agent{\Hyp_t}(\spec{C}) \ \text{expandable by scaffold rewrites};
\qquad
\text{(M2)}\ \ \val{V}(\tgt{z}) \le \val{V^*(\agent{W})}\ \text{for all reachable } \tgt{z}.
\]
\[
\sum_{t=1}^{\infty}\val{\eta_t} \ \le\ \val{V^*(\agent{W})} - \mathbb{E}[\val{V}(\tgt{z_0})],
\qquad \mathbb{E}[\val{V}(\tgt{z_t})] \nearrow \val{L} \le \val{V^*(\agent{W})}.
\]
\begin{legend}
\leg{$\agent{\Hyp_t}(\spec{C})$}{effective hypothesis class of reachable patches, a function of the mutable scaffold $\spec{C}$.}
\leg{$\agent{W}$}{the model weights (frozen).}
\leg{$\val{V^*(\agent{W})}$}{the ultimate quality ceiling achievable with weights $\agent{W}$ by \emph{any} scaffold.}
\leg{$\val{L}$}{the limit the system saturates to — no greater than the ceiling.}
\end{legend}
\xref{Corollary~4 (Frozen-Weight Ceiling): frozen weights guarantee the ceiling $\val{V^*(\agent{W})}$, not a safe trajectory to it. Lean: \texttt{cor\_frozen\_weight\_ceiling}.}

\section{Orchestration Selects a Ceiling; It Does Not Raise One}

\[
\val{V^*_{\mathrm{orch}}} \;=\; \max_i \val{V^*(\agent{W},\spec{C_i})}
\qquad\text{(exact equality, unconditional)}.
\]
\[
\mathbb{E}\big[\val{V}(\tgt{z_T})\big] \ \ge\ \val{V}(\tgt{z_0}) + \sum_{t=1}^{T}\max_i \val{\eta_t^{(i)}},
\qquad
\sum_{t=1}^{T}\max_i \val{\eta_t^{(i)}} \ \le\ 1 - \val{V}(\tgt{z_0}).
\]
\begin{legend}
\leg{$\spec{C_i}$}{worker $i$'s scaffold — distinct prompts, tools, or specializations, possibly sharing weights $\agent{W}$.}
\leg{$\val{V^*(\agent{W},\spec{C_i})}$}{worker $i$'s tight ceiling, $\sup\{\val{V}(\tgt{z}):\tgt{z}$ reachable under $\spec{C_i}\}$.}
\leg{$\val{V^*_{\mathrm{orch}}}$}{what best-of-$k$ selection attains. The $=$ is the content: orchestration \emph{realizes} the best worker's ceiling and never exceeds it, even when no single worker's reachable class contains the union.}
\leg{$\max_i \val{\eta_t^{(i)}}$}{the orchestrator banks each round's best gain, dominating every worker's depth-only guarantee\ldots}
\leg{$\le 1-\val{V}(\tgt{z_0})$}{\ldots yet the total is still capped by the same bounded metric: \textbf{width buys rate, not budget}.}
\end{legend}
\xref{Proposition~5 (Orchestration Selects a Ceiling). The equality is where the formalization
earned its place: a conjectured \emph{strict} inequality was refuted by two disjoint singleton
classes of equal value. Width half stated separately as Prop.~18, whose extra hypothesis
(monotonicity of the inner action-expectation) the formalization flags.}

\section{Bounded Overlap Characterizes a Correct Vote}

\[
\erf{m_i} \;:=\; \#\{\, 0\le t<k \;:\; \tgt{y_i}\,\agent{h_t}(\spec{x_i}) = -1 \,\},
\qquad
\erf{m^\star} \;:=\; \max_i \erf{m_i}.
\]
\[
\tgt{y_i}\,\wt{f}(\spec{x_i}) \;=\; k - 2\erf{m_i}
\qquad\Longrightarrow\qquad
\text{vote error}=0
\ \iff\ 2\erf{m_i}<k\ \ \forall i
\ \iff\ \erf{m^\star} < k/2 .
\]
\[
\text{measured (Part~B): } \erf{m^\star} = k \ \text{ at every round.}
\]
\begin{legend}
\leg{$\erf{m_i}$}{how many of the $k$ workers fail point $i$. The index runs $0\le t<k$, matching the formalization's \texttt{range k}; it counts all $k$ workers, not $k-1$.}
\leg{$\erf{m^\star}$}{the \emph{tight overlap} — one integer that decides whether a vote can work at all.}
\leg{$\tgt{y_i}\wt{f}(\spec{x_i})=k-2\erf{m_i}$}{splitting the unweighted margin at point $i$; this identity is the whole proof.}
\leg{$\erf{m^\star}<k/2$}{\emph{exact}, not merely sufficient — and sharp: one point at $2\erf{m_i}{=}k$ already makes the error nonzero, and if every point sits there the error is $1$.}
\leg{$\erf{m^\star}=k$}{our pool at the ceiling of overlap: some task defeats \emph{every} worker, so the red mass is the vote's error.}
\end{legend}
\xref{Proposition~6 (Bounded Overlap Characterizes a Correct Vote), measured in
\S\,Theoretical Analysis (Fig.~1). Weighted voting is the one boosting-\emph{exact}
arrangement (Prop.~17), which is why this condition, not the ceiling, is what binds a
verifier-free orchestrator. Lean: \texttt{vote\_margin\_eq},
\texttt{thm\_overlap\_iff\_zero\_error}, \texttt{cor\_tight\_overlap\_iff},
\texttt{thm\_overlap\_boundary\_sharp}.}

\section{The Reused-Pool Dichotomy}

\[
\erf{\epsilon_0} \;:=\; \min\Big\{\text{zero-one training error of } \textstyle\sum_j \wt{A_j}\,\agent{H_j} \;:\; \wt{A_j}\ge 0\Big\}
\qquad\text{(a property of the pool alone)}.
\]
\[
\text{(i)}\ \ \erf{\epsilon_0}>0:\quad
\erf{\epsilon_0} \ \le\ e^{-2\erf{\gamma}^2 T}
\qquad\Longrightarrow\qquad
T \ \le\ \frac{\log(1/\erf{\epsilon_0})}{2\erf{\gamma}^2}.
\]
\[
\text{(ii)}\ \ \erf{\epsilon_0}=0:\quad
\forall \spec{D}\ \ \exists j:\ \sum_{i\,:\,\tgt{y_i}\agent{H_j}(\spec{x_i})=-1}\!\!\spec{D}(i) \ \le\ \tfrac12-\erf{\gamma},
\qquad \erf{\gamma} \;=\; \frac{\erf{\theta}}{2\wt{S}} \;>\; 0 .
\]
\begin{legend}
\leg{$\erf{\epsilon_0}$}{best zero-one training error any \emph{nonnegative combination} of the pool attains — the quantity that decides which branch applies.}
\leg{(i)}{the \emph{horizon}: the edge hypothesis provably fails by a finite, computable round, uniformly over the schedule and the confidences, depending on $k$ only through $\erf{\epsilon_0}$.}
\leg{$\erf{\theta},\wt{S}$}{interpolation margin $\min_i \tgt{y_i}\sum_j \wt{A_j}\agent{H_j}(\spec{x_i})$ and coefficient sum $\sum_j \wt{A_j}$ of the witnessing combination.}
\leg{(ii)}{the edge $\erf{\gamma}{=}\erf{\theta}/2\wt{S}$ is fixed \emph{before} the distribution is chosen, so the hypothesis survives every round of any run.}
\leg{exhaustive}{the two cases are mutually exclusive and cover everything: \textbf{a fixed pool is not a renewable resource}.}
\end{legend}
\xref{Proposition~7 (Reused-Pool Dichotomy); halves stated with their exact hypotheses as
Props.~20 and~21. Only the interpolating direction is formalized — the converse needs
AdaBoost's minimax duality and remains unformalized.}

\section{Framework Mapping (Boosting $\leftrightarrow$ Coding)}

\[
\underbrace{\spec{X}}_{\text{input}}\ \mapsto\ \underbrace{\tgt{Y}}_{\text{target}},
\qquad
\underbrace{\tgt{Y}-\agent{F_{t-1}}(\spec{X})}_{\text{residual}}\ \approx\ \agent{h_t},
\qquad
\text{loss } \erf{L}(\tgt{Y},\hat{\tgt{Y}}),\quad
\text{distribution } \spec{D_t}.
\]
\begin{legend}
\leg{$\spec{X}\mapsto\tgt{Y}$}{structured-prediction view: a specification maps to its optimal code.}
\leg{$\agent{F_{t-1}}(\spec{X})$}{current draft after $t-1$ rounds.}
\leg{$\tgt{Y}-\agent{F_{t-1}}(\spec{X})$}{the \emph{residual error} each weak agent $\agent{h_t}$ predicts (a git diff / patch).}
\end{legend}
\xref{Table~1 (Framework mapping), \S\,Framework.}

\section{Algorithm: Agentic Boosting Protocol for Specifications}

\begin{align*}
&\textbf{Require: } \{(\spec{x_i},\tgt{y_i})\}_{i=1}^{N},\ T
   &&\text{spec--code pairs, rounds}\\
&\spec{w_i} \leftarrow 1/N
   &&\text{initialize weights}\\
&\erf{\epsilon_t} = \sum_i \spec{w_i}\,\mathbf{1}\!\big[\erf{L}(\tgt{y_i},\agent{h_t}(\spec{x_i}))>0\big]
   &&\text{select } \agent{h_t}\text{ minimizing weighted error}\\
&\text{if } \erf{\epsilon_t}\ge \tfrac12:\ \text{stop}
   &&\text{no weak learner available}\\
&\wt{\alpha_t} \leftarrow \tfrac12\ln\frac{1-\erf{\epsilon_t}}{\erf{\epsilon_t}}
   &&\text{confidence weight}\\
&\spec{w_i} \leftarrow \spec{w_i}\cdot\exp\!\big(-\wt{\alpha_t}\,\mathbf{1}[\erf{L}(\tgt{y_i},\agent{h_t}(\spec{x_i}))\,{=}\,0]\big)
   &&\text{re-weight (note: }{=}0\text{ = success)}\\
&\spec{w_i} \leftarrow \spec{w_i}/\textstyle\sum_j \spec{w_j}
   &&\text{normalize}\\
&\textbf{return } \agent{H}(\spec{x})=(\agent{R_T}\circ\cdots\circ \agent{R_1})(\tgt{z_0}),\ \ \wt{f}(\spec{x})=\textstyle\sum_t \wt{\alpha_t}\,\agent{h_t}(\spec{x})
   &&\text{artifact + margin}
\end{align*}
\begin{legend}
\leg{$\spec{w_i}$}{unnormalized weight on spec $i$ (the operational form of $\spec{D_t}$).}
\leg{$\mathbf{1}[\erf{L}=0]$}{success indicator — note this flips the theory's $>0$ (failure) convention.}
\leg{return}{artifact $\agent{H}$ by sequential composition; confidence by additive margin $\wt{f}$.}
\end{legend}
\xref{Algorithm~1 (Appendix). Operational form of Definition~9.}

\section{Part~B: SWE-bench Regime Estimators}

\[
\val{q_{i,t}}\in[0,1],\quad \erf{\tau}=0.6,\qquad
\erf{\err_t}=\tfrac{1}{K}\sum_i \mathbf{1}\!\big[\max_{s\le t}\val{q_{i,s}}<\erf{\tau}\big]\ \approx\ \exp\!\big(-2\textstyle\sum_s \erf{\gamma_s}^2\big).
\]
\[
\val{S_{i,t}}=\sum_{s\le t}\max\!\big(0,\val{q_{i,s}}-\val{q_{i,s-1}}\big)\ \le\ 1-\val{q_{i,0}},
\qquad
\val{m_{i,t}}=2\val{q_{i,t}}-1.
\]
\begin{legend}
\leg{$\val{q_{i,t}}$}{binary harness pass/fail of task $i$ at round $t$ ($1$ = all \texttt{FAIL\_TO\_PASS} pass, no regressions).}
\leg{$\erf{\tau}=0.6$}{pass \emph{threshold} on the $[0,1]$ partial score $\val{q_{i,t}}$ — the only meaning $\tau$ now carries.}
\leg{$\erf{\err_t}$}{running-best pass error over $K$ tasks; tracks the exponential-decay prediction.}
\leg{$\val{S_{i,t}}$}{cumulative improvement, bounded by headroom $1-\val{q_{i,0}}$ (refinement saturation).}
\leg{$\val{m_{i,t}}$}{margin per task; predicted to shift rightward across rounds.}
\end{legend}
\xref{\S\,Part~B (Quality signal; Predictions), $K{=}\statPBK$, $T{=}\statPBRounds$.}

\section{Part~C: Production-Session Empirical Model}

\[
\text{survivor ratio} = \frac{|\text{net lines}|}{\text{total churn}},
\qquad
\empc{\psi_{\mathrm{Opus}}}\approx \statPCTauOpus \ \ \text{vs.}\ \ \empc{\psi_{\mathrm{Sonnet}}}\approx \statPCTauSonnet.
\]
\[
\log(\text{churn}_{t+1}) \sim \mathcal{N}\!\big(\empc{\rho}\log(\text{churn}_t),\,\empc{\sigma}^2\big),
\quad \empc{\rho}=\statPCRho,\ \empc{\sigma}=\statPCSigma,\quad
\empc{\pi_{\mathrm{pure}}}\approx \statPCPiPure.
\]
\begin{legend}
\leg{survivor ratio}{structural proxy: net lines retained over total churn (\texttt{git\_session\_extractor.py:152}).}
\leg{$\empc{\psi_{\mathrm{Opus}}},\empc{\psi_{\mathrm{Sonnet}}}$}{mean non-pure survivor ratios by model tier — the Part-C meaning of $\tau$ (a statistic, not a threshold). Gap $=\statPCGap$ (95\%~CI \statPCGapCI), a null surviving four independent robustness checks --- not a directional effect.}
\leg{$\empc{\rho}$}{geometric churn-decay rate (log-normal AR(1) coefficient); $<1$ signals convergence.}
\leg{$\empc{\sigma}$}{spread of the log-normal churn model.}
\leg{$\empc{\pi_{\mathrm{pure}}}$}{mixture weight for pure-addition sessions (survivor ratio $=1$).}
\end{legend}
\xref{\S\,Part~C, Findings~1--3. Ties $\empc{\rho}$ to Theorem~2 decay and $\empc{\psi}$ to the $\val{V^*(\agent{W})}$ ceiling (Corollary~4).}

\end{document}
